Streamers 已在多個正在形成的原行星盤周圍，透過高解析度的 ALMA 觀測被發現。先前研究指出，Streamers 可能在自坍縮的前恆星核心向內吸積物質的過程中，主導整體的質量吸積收支，並對原行星盤的動力學結構產生重要影響。因此，理解 Streamers 的形成機制，對於解釋原行星盤如何從包層獲得質量具有關鍵意義。我們提出一個模型，透過考慮由重力不穩定所導致的密度增強，來解釋 Streamer 結構的形成。我們將此模型應用於已觀測到 Streamers 的天體進行測試，包括以 NOEMA 觀測的 Per-emb-2 與 Per-emb-50，以及以 ALMA 觀測的 SCrA 與 HLTau，並對模型參數進行擬合。此方法使我們能更深入探討 Streamers 的物理起源，以及它們在將質量由核心輸送至原行星盤過程中所扮演的角色。該模型亦可進一步應用於分析新的觀測資料，以及具有 Streamer 特徵的既有觀測資料庫。